\subsection{Pagine di errore}

In alcune situazioni particolari, l'applicazione pu\`o mostrarti una pagina di errore.\\
Queste pagine servono a spiegarti cosa \`e successo e a indicarti come proseguire senza problemi.

\subsubsection{Pagina 404 -- Pagina non trovata}

La pagina 404 viene mostrata quando provi ad accedere a una pagina che non esiste o non \`e pi\`u disponibile.

\begin{figure}[H]
	\centering
	\includegraphics[width=0.5\textwidth]{ManualeUtente/errore_404}
	\caption{Pagina 404 - Pagina non trovata}
\end{figure}

In pratica, significa che il sistema non riesce a trovare la risorsa richiesta.\\
Questo pu\`o succedere, ad esempio, se:
\begin{itemize}[itemsep=5pt, parsep=5pt, label=$\scriptstyle\bullet$]
	\item L'indirizzo inserito non \`e corretto
	\item La pagina non \`e pi\`u disponibile
	\item Hai seguito un link non valido
\end{itemize}

Per continuare, puoi scegliere una delle seguenti opzioni:

\begin{itemize}[itemsep=5pt, parsep=5pt, label=$\scriptstyle\bullet$]
	\item \textbf{Tornare alla \hyperlink{glossario:chat}{chat}}: ti riporta alla schermata principale dell'applicazione.
	\item \textbf{Vai al \hyperlink{glossario:login}{login}}: ti permette di effettuare nuovamente l'accesso al sistema.
\end{itemize}

Queste opzioni ti aiutano a riprendere rapidamente la navigazione.

\subsubsection{Pagina 401 -- Non autorizzato}

La pagina 401 viene mostrata quando provi ad accedere a una sezione per cui non hai i permessi necessari.

\begin{figure}[H]
	\centering
	\includegraphics[width=0.5\textwidth]{ManualeUtente/errore_401}
	\caption{Pagina 401 - Non autorizzato}
\end{figure}

Questo pu\`o accadere, ad esempio, se:
\begin{itemize}[itemsep=5pt, parsep=5pt, label=$\scriptstyle\bullet$]
	\item Non hai effettuato l'accesso
	\item La tua \hyperlink{glossario:sessione}{sessione} \`e scaduta
	\item Stai tentando di accedere a una sezione riservata, come l'area \hyperlink{glossario:amministratore}{amministratore}
\end{itemize}

In questo caso, puoi:

\begin{itemize}[itemsep=5pt, parsep=5pt, label=$\scriptstyle\bullet$]
	\item \textbf{Tornare al \hyperlink{glossario:login}{login}}: ti permette di accedere nuovamente al sistema con le tue credenziali.
\end{itemize}
